\documentclass{article}
\usepackage{graphicx} % Required for inserting images
\usepackage{caption} % Required for \captionof
\usepackage{amsmath}
\usepackage{float}
\usepackage{comment}
\usepackage[T1]{fontenc}
\usepackage[utf8]{inputenc}
\usepackage{lmodern}
\usepackage{microtype}
\usepackage{xcolor}
\usepackage{framed}
\renewenvironment{leftbar}{%
  \def\FrameCommand{\textcolor{orange}{\vrule width 3pt} \hspace{10pt}}%
  \MakeFramed {\advance\hsize-\width \FrameRestore}}%
 {\endMakeFramed}

\title{Mechanikai mérés}
\author{Kern Luca, Vincze Csongor}
\date{2026 Április 22}

\begin{document}

\maketitle

\section*{7. feladat: Határozza meg a rendszer $\omega_0$ saját körfrekvenciáját a csillapítatlan szabad rezgés periódusidejének mérésével!}

A csillapító mágneseket továbbra sem szereltük vissza, az ultrahangos távolságmérőt rögzítettük. A rudat úgy rögzítettük, hogy a közepe kb a rudvezetőnél van, és a rúdvezetővel nem érintkezik. Rézlapot ragasztottunk a mérőrúd aljára, hogy növeljük a felületet a távolságméréshez.
A rudat $x = \text{mm}$-rel kitérítettük, és megmértük a távolságot az időfüggvényében 5-10 perióduson keresztül.
$m_{rúd} = 50 \text{g}$

\begin{table}[H]
    \centering
    \renewcommand{\arraystretch}{1.2}
    \begin{tabular}{|l|c|c|c|}
        \hline
        Mérés száma & Kitérés [mm] & Idő [s] & $\omega_0$ [rad/s] \\ \hline
        1 & & & \\ \hline
        2 & & & \\ \hline
        3 & & & \\ \hline
    \end{tabular}
    \caption{A rendszer saját körfrekvenciájának meghatározása}
    \label{tab:mecha_7}
\end{table}

Ugyanezt megismétetltük $m = 50\text{g}$ tömeggel a rúdra helyezve.

\begin{table}[H]
    \centering
    \renewcommand{\arraystretch}{1.2}
    \begin{tabular}{|l|c|c|c|}
        \hline
        Tömeg [g] & Kitérés [mm] & Idő [s] & $\omega_0$ [rad/s] \\ \hline
        50 & & & \\ \hline
        100 & & & \\ \hline
        150 & & & \\ \hline
    \end{tabular}
    \caption{A rendszer saját körfrekvenciájának meghatározása tömeggel}
    \label{tab:mecha_7_2}
\end{table}

\end{document}